% ============================================================================
% Fichier  : partie6/P06_C26_mesurer.tex
% Titre    : Mesurer la Sécurité Offensive
% Auteur   : MINKA MI NGUIDJOÏ Thierry Emmanuel
% Version  : 1.0 -- novembre 2026
% ============================================================================
% Sources : cadre de verdict OUI / NON / NEK de l'auditologie (MINKA) ;
% indicateurs de sécurité offensive. Opérations LOGONE et MBAM : cas fictifs.
% ============================================================================

\chapter{Mesurer la Sécurité Offensive}
\label{chap:mesurer}
\chaptermark{Mesurer la sécurité offensive}

\epigraph{%
    « Un test qui ne se mesure pas ne s'améliore pas~;\\
    et ce qui ne peut être tranché doit être dit indéterminé. »%
}{Principe du cours}

\niveauChapitre{Avancé}{L'ensemble du
cours~; notions d'auditologie (cours d'audit).}

\begin{boxobjectifs}
\begin{itemize}
    \item Définir ce que mesure un test d'intrusion.
    \item Construire des indicateurs utiles et honnêtes.
    \item Appliquer le verdict à trois valeurs \sigle{OUI}, \sigle{NON},
          \sigle{NEK}.
    \item Éviter les indicateurs trompeurs.
    \item Inscrire le test dans un cycle d'amélioration.
\end{itemize}
\end{boxobjectifs}

\bigskip

% ============================================================================
\section*{Question directrice : LiZbiZ est-elle « sûre » après le test~?}
% ============================================================================

À la fin de la mission, M.~ETEME pose la question la plus simple et la plus
piégée~: « sommes-nous sûrs maintenant~? ». La réponse honnête n'est ni oui
ni non en général~: elle est un ensemble de verdicts précis, et, pour ce qui
n'a pu être tranché, un franc « indéterminé ». Ce chapitre apprend à
mesurer sans tromper.

\medskip
\textbf{Plan du chapitre.}
\begin{enumerate}
    \item Ce que mesure un test (section~\ref{sec:c26-quoi}).
    \item Des indicateurs honnêtes (section~\ref{sec:c26-indicateurs}).
    \item Le verdict à trois valeurs (section~\ref{sec:c26-verdict}).
    \item Les indicateurs trompeurs (section~\ref{sec:c26-trompeurs}).
    \item Le cycle d'amélioration (section~\ref{sec:c26-cycle}).
\end{enumerate}

\separateur

% ============================================================================
\section{Ce que Mesure un Test}
\label{sec:c26-quoi}
% ============================================================================

Un test d'intrusion ne mesure pas « la sécurité » en général~: il mesure la
résistance d'un périmètre donné, à un moment donné, face à un effort donné.
Il dit ce qu'un attaquant de niveau comparable aurait pu faire, pas ce qu'un
attaquant quelconque fera demain.

% ============================================================================
\section{Des Indicateurs Honnêtes}
\label{sec:c26-indicateurs}
% ============================================================================

\begin{tableaucours}{L{4.8cm}L{7.4cm}}{%
    \tableauHeader{Indicateur} &
    \tableauHeader{Ce qu'il dit}
}
Objectifs atteints sur objectifs fixés & Résistance aux scénarios visés \\
Constats par gravité & Ampleur et priorité des faiblesses \\
Couverture de détection & Part des techniques vues par le défenseur \\
Temps jusqu'à la détection & Rapidité du défenseur \\
Taux de correction au contre-test & Qualité du traitement \\
Taux de signalement d'hameçonnage & Vigilance humaine \\
\end{tableaucours}
\vspace{-0.8em}
\begin{center}{\small\itshape Tableau~26.1. Indicateurs d'un test d'intrusion}\end{center}

% ============================================================================
\section{Le Verdict à Trois Valeurs}
\label{sec:c26-verdict}
% ============================================================================

\begin{boxdef}
\termecle{Verdict à trois valeurs}~: pour chaque question d'audit, la réponse
est \sigle{OUI} (la propriété est vérifiée, avec preuve), \sigle{NON} (elle
est réfutée, avec preuve), ou \sigle{NEK} (indéterminé~: les preuves
manquent pour trancher). Le \sigle{NEK} n'est pas un aveu de faiblesse~:
c'est une exigence d'honnêteté.
\end{boxdef}

\begin{tableaucours}{L{2.0cm}L{4.6cm}L{5.6cm}}{%
    \tableauHeader{Verdict} &
    \tableauHeader{Sens} &
    \tableauHeader{Exemple (LOGONE)}
}
OUI & Propriété démontrée & « Le système de paiement cloisonné a résisté »
(F3 non atteint, MBAM) \\
NON & Propriété réfutée & « Un attaquant interne devient administrateur du
domaine » (T-01) \\
NEK & Indéterminé & « La résistance du nuage au-delà du périmètre testé » (hors
mandat) \\
\end{tableaucours}
\vspace{-0.8em}
\begin{center}{\small\itshape Tableau~26.2. Le verdict OUI / NON / NEK appliqué~\cite{minka2026auditologie}}\end{center}

\begin{boxalert}
\textbf{Violation typique.} Conclure \sigle{OUI} (« c'est sûr ») faute
d'avoir trouvé une faille, alors qu'on n'a pas cherché.

\smallskip
\textbf{Règle~:} l'absence de preuve n'est pas une preuve d'absence. Ce qui
n'a pas été testé, ou n'a pu être tranché, est \sigle{NEK}, pas \sigle{OUI}.
Le \sigle{NEK} dit au client où porter l'effort suivant.
\end{boxalert}

% ============================================================================
\section{Les Indicateurs Trompeurs}
\label{sec:c26-trompeurs}
% ============================================================================

\begin{tableaucours}{L{4.4cm}L{7.8cm}}{%
    \tableauHeader{Indicateur trompeur} &
    \tableauHeader{Pourquoi il trompe}
}
Nombre brut de failles & Dépend de l'outil, pas du risque \\
« Zéro faille trouvée » & Peut signifier « rien cherché » \\
Note globale unique & Masque les constats critiques dans la moyenne \\
Temps passé & Mesure l'effort, pas le résultat \\
\end{tableaucours}
\vspace{-0.8em}
\begin{center}{\small\itshape Tableau~26.3. Indicateurs à manier avec prudence}\end{center}

% ============================================================================
\section{Le Cycle d'Amélioration}
\label{sec:c26-cycle}
% ============================================================================

Un test n'est pas une fin~: c'est une étape d'un cycle. Les constats
alimentent les corrections~; le contre-test les vérifie~; la phase violette
nourrit la détection~; le test suivant repart des \sigle{NEK} et des zones
non couvertes. La sécurité offensive sert la défense, et se mesure à ce
qu'elle lui apporte.

\begin{boxmetier}{RSSI, auditeur de sécurité, pentesteur}
\textbf{Dire le vrai, y compris le NEK.} Le professionnel ne cède pas à la
pression du \sigle{OUI} rassurant. Il distingue ce qu'il a démontré, ce qu'il
a réfuté, et ce qu'il n'a pu trancher~: c'est à ce triple langage qu'on
reconnaît un rapport honnête.
\end{boxmetier}

\separateur

\begin{boxresume}
\textbf{Ce qu'il faut retenir de ce chapitre~:}
\begin{itemize}
    \item \textbf{Un test} mesure une résistance, à un moment, face à un
          effort donné.
    \item \textbf{Les bons indicateurs} relient objectif, gravité, détection,
          correction.
    \item \textbf{Le verdict} est \sigle{OUI}, \sigle{NON} ou \sigle{NEK}.
    \item \textbf{L'absence de preuve} n'est pas preuve d'absence~: c'est
          \sigle{NEK}.
    \item \textbf{Le test} sert un cycle d'amélioration continue de la
          défense.
\end{itemize}
\end{boxresume}

\begin{boxinfo}
\textbf{Points de contrôle de l'élève.} Je sais~: dire ce que mesure un
test~; construire des indicateurs honnêtes~; appliquer OUI / NON / NEK~;
reconnaître les indicateurs trompeurs~; inscrire le test dans un cycle.
\end{boxinfo}

\bigskip

\begin{boxexo}[title={\ding{45}\quad Exercice~26.1. Ce que mesure un test \niveaubase}]
Un test mesure-t-il « la sécurité » en général~? Que mesure-t-il~?
\end{boxexo}

\begin{boxexo}[title={\ding{45}\quad Exercice~26.2. Trois valeurs \niveaubase}]
Donnez le sens de OUI, NON et NEK.
\end{boxexo}

\begin{boxexo}[title={\ding{45}\quad Exercice~26.3. Trompeur \niveaubase}]
Pourquoi « zéro faille trouvée » peut-il tromper~?
\end{boxexo}

\begin{boxexo}[title={\ding{45}\quad Exercice~26.4. Absence de preuve \niveaumoyen}]
Expliquez « l'absence de preuve n'est pas une preuve d'absence » sur un
exemple de LOGONE.
\end{boxexo}

\begin{boxexo}[title={\ding{45}\quad Exercice~26.5. Classer \niveaumoyen}]
Attribuez OUI, NON ou NEK~: le système de paiement de MBAM~; l'accès
administrateur de LiZbiZ~; la sécurité du nuage hors périmètre.
\end{boxexo}

\begin{boxexo}[title={\ding{45}\quad Exercice~26.6. Note unique \niveaumoyen}]
Pourquoi une note globale unique est-elle dangereuse~?
\end{boxexo}

\begin{boxexo}[title={\ding{45}\quad Exercice~26.7. Indicateurs \niveaumoyen}]
Choisissez trois indicateurs pour le tableau de bord de LiZbiZ et justifiez.
\end{boxexo}

\begin{boxexo}[title={\ding{45}\quad Exercice~26.8. Répondre au DG \niveaumoyen}]
Rédigez la réponse honnête à « sommes-nous sûrs maintenant~? », en OUI / NON
/ NEK.
\end{boxexo}

\begin{boxexo}[title={\ding{45}\quad Exercice~26.9. Cycle \niveauexpert}]
Décrivez le cycle d'amélioration reliant test, correction, contre-test,
détection et test suivant.
\end{boxexo}

\begin{boxexo}[title={\ding{45}\quad Exercice~26.10. Tableau de bord \niveauexpert}]
Construisez le tableau de bord de sécurité offensive que LiZbiZ devrait
suivre d'une année sur l'autre, verdicts NEK compris.
\end{boxexo}

\bigskip

\begin{boxinfo}
\textbf{Pour aller plus loin.}
\begin{itemize}
    \item Cadre d'auditologie et verdict à trois valeurs~\cite{minka2026auditologie}.
\end{itemize}
\end{boxinfo}

\begin{boxinfo}
\textbf{Fin du cours.} De la lettre d'autorisation au verdict mesuré, le
test d'intrusion aura montré qu'autoriser, reconnaître, tester et restituer
forment un même service~: défendre, en pensant comme l'attaquant, sans
jamais le devenir.
\end{boxinfo}
